\section{Introduction}
\label{paper:introduction}

For real numbers $3\leq N_1\leq\cdots\leq N_{k+1}$, write
\[
 A_{k+1}(\mathbf N)
 =\#\{n_1\cdots n_{k+1}: n_i\in\mathbb N,\ n_i\leq N_i\}.
\]
The generalized multiplication-table problem asks for the order of this
quantity, uniformly in the side lengths, with the number of factors fixed.
Although there are about $N_1\cdots N_{k+1}$ ways to select the factors,
many selections have the same product. The problem is to measure this
multiplicity without losing the dependence on the relative sizes of the
sides. We determine the order of $A_{k+1}(\mathbf N)$ for every fixed
$k\geq6$ and all ordered side lengths. Combined with the work of Ford and
Koukoulopoulos, this completes the order-of-magnitude theory in every
fixed dimension.

\subsection{The arithmetic results}

The divisor-interval formulation both gives a more general arithmetic
result and reveals the structure of the answer. For
$\mathbf y=(y_1,\ldots,y_k)$, let
\[
 H^{(k+1)}(x,\mathbf y,2\mathbf y)
 =\#\bigl\{n\in\mathbb N:n\leq x,\
 \exists d_i\in\mathbb N\cap(y_i,2y_i]\ (1\leq i\leq k),\
 d_1\cdots d_k\mid n\bigr\}.
\]
All logarithms in this paper are natural. Set $y_0=3$ and define
\begin{equation}
 \ell_i(\mathbf y)=\log\frac{3\log y_i}{\log y_{i-1}},
 \qquad
 L_i=\ell_i-\log3=\log\frac{\log y_i}{\log y_{i-1}}.
 \label{paper:ell}
\end{equation}
Thus ordered cutoffs $3\leq y_1\leq\cdots\leq y_k$ correspond exactly
to vectors $\ell\in[\log3,\infty)^k$, or equivalently $L_i\geq0$.

The answer is expressed by a single function $\Phi_k$ of these logarithmic
lengths. Its definition, given in Section~\ref{profile:definition} and
completed in \eqref{profile:final}, involves a finite-dimensional ordered
optimization, finitely many endpoint comparisons, and a scalar
conditional persistence exponent. In particular, $\Phi_k$ is defined
independently of the arithmetic quantities being estimated. We first
state the two arithmetic consequences; the ingredients of the profile
are explained below and then specified in full.

\begin{theorem}[Divisors in prescribed dyadic intervals]
\label{paper:main-H}
Fix an integer $k\geq6$. Uniformly for real parameters satisfying
\begin{equation}
 3\leq y_1\leq\cdots\leq y_k,
 \qquad
 x\geq2^k\left(\prod_{i=1}^k y_i\right)y_k,
 \label{paper:ADH}
\end{equation}
one has
\[
 \frac{H^{(k+1)}(x,\mathbf y,2\mathbf y)}{x}
 \asymp_k \Phi_k\bigl(\ell(\mathbf y)\bigr).
\]
\end{theorem}

\begin{theorem}[Rectangular multiplication tables]
\label{paper:main-table}
Fix an integer $k\geq6$. For all real side lengths
$3\leq N_1\leq\cdots\leq N_{k+1}$, put
\[
 Y_0=3,\qquad Y_i=\max\{3,N_i/2\}\quad(1\leq i\leq k).
\]
Then
\[
 A_{k+1}(\mathbf N)
 \asymp_k
 \left(\prod_{i=1}^{k+1}N_i\right)
 \Phi_k\bigl(\ell(\mathbf Y)\bigr).
\]
\end{theorem}

Here $f\asymp_k g$ means that $c_kg\leq f\leq C_kg$ for positive
constants depending only on $k$. In particular, repeated side lengths,
bounded initial sides and arbitrarily large ratios of sides are all
included. The dimension of the table is $k+1$, so the new range starts
with seven factors. The constants are allowed to depend on the fixed
dimension. The restriction on $x$ in \eqref{paper:ADH} is the precise
range of the divisor-count comparison used in the proof;
Theorem~\ref{paper:main-table} covers the full stated range of side
lengths, including those for which some $N_i/2<3$.

\subsection{Previous work and the remaining dimensions}

The classical problem goes back to Erd\H{o}s~\cite{Erdos1955}.
Its connection with the distribution of divisors in an interval led to
extensive work on $H^{(2)}(x,y,z)$, including Tenenbaum's
estimates~\cite{Tenenbaum1984}. Ford determined the order of the
classical square table~\cite[Corollary~3]{FordDivisors2008}. With
$Q(u)=u\log u-u+1$, his result is
\[
 A_2(N,N)\asymp
 \frac{N^2}{(\log N)^{Q(1/\log2)}(\log\log N)^{3/2}}
 \qquad(N\geq3).
\]
Ford's interval estimate~\cite[Corollary~2]{FordDivisors2008}, together
with the multiplication-table comparison in
\cite[Theorem~1.1]{Kouk}, also gives the order of all two-dimensional
rectangular tables. This distinction between squares and rectangles is
useful in describing the higher-dimensional problem.

Koukoulopoulos proved that for every fixed $k\geq2$,
\begin{equation}
 A_{k+1}(N,\ldots,N)\asymp_k
 \frac{N^{k+1}}
 {(\log N)^{Q(k/\log(k+1))}(\log\log N)^{3/2}},
 \qquad N\geq3,
 \label{paper:known-cube}
\end{equation}
see \cite[Corollary~1]{KoukLocal}. His corresponding divisor theorem
covers comparable logarithmic scales. Thus high-dimensional cubes were
already understood. For arbitrary side lengths, his later work
\cite{Kouk} established a uniform estimate throughout
\eqref{paper:ADH} when $2\leq k\leq5$, and consequently determined the
order of rectangular tables with three through six factors.
The same paper gives a general upper bound in every fixed dimension,
a matching estimate under an additional condition on the saddle
parameter, and a high-dimensional range in which the logarithmic scales
are sufficiently close; see its Theorems~1.3--1.5. Its
Section~2.2, Lemma~2.3 and Remark~2.4 identify finer structural
constraints that can make the earlier expression too large when
$k\geq6$. The task of determining the uniform order for these remaining
dimensions was also stated explicitly in
\cite[Chapter~6, pp.~87--88]{KoukThesis2010}.
Theorems~\ref{paper:main-H} and~\ref{paper:main-table} resolve these
remaining fixed-dimensional cases.

More precise questions about the square table have recently seen
further progress. Sawhney announced joint work with Green giving an
asymptotic formula with a periodic amplitude for the two-dimensional
integer square table~\cite{SawhneyGreenAnnouncement2025}. Their
paper~\cite{GreenSawhney2026} proves an asymptotic formula for the
permutation analogue and describes the integer-table work as
forthcoming. Sabuncu and T\'afula~\cite{SabuncuTafula2026} study the
regime in which the dimension grows sufficiently rapidly relative to
a common side length. The present results concern uniform two-sided
bounds for arbitrary side lengths in each fixed dimension.

\subsection{What the profile records}

The exponential part of the answer is governed by an ordered saddle
$s_1\geq\cdots\geq s_k$. Its maximal consecutive blocks of equal
coordinates organize the constraints; we call these blocks
\emph{owners}. Equality of saddle coordinates does not identify the
original prime ranges or the divisor coordinates. Those ranges retain
their own lengths, and the coordinates retain their own deadlines.
This information is needed to distinguish several losses that can occur
at the same exponential scale.

For positive lengths, the profile has the form
\[
 \Psi_k(L)=e^{-J(L)}B_{\mathrm{ord}}
 \prod_{\substack{\text{proper}\\\text{owners}}}
       (1+d\sqrt{1+T})^{-\sigma}
 \prod_{\text{canonical edges }e}P_e.
\]
The saddle cost $J$ supplies the principal exponential rate. A proper
owner, meaning one that does not contain the last coordinate, contributes
a count-window factor: $T$ is its total native length, $\sigma$ its
saddle value, and $d$ its relative gap from the next owner. This factor
keeps the transition between a fixed gap and a gap of order $T^{-1/2}$
in the same formula. The terminal owner contributes one ordinary factor
$B_{\mathrm{ord}}$, which includes the terminal count probability and
its order-statistics bridge.

The remaining factors describe binary comparisons between a set of
retained names and its complement. Their endpoint scales are obtained
from a finite enumeration of legal histories on the original groups.
The associated paths share a random environment: one first fixes the
common gaps, averages over the private marks, and then takes a
fractional moment. This order of averaging produces the exponent
$\chi(\eta,\rho)$ in \eqref{profile:chi}, defined using two independent
stationary Ornstein--Uhlenbeck processes. The parameters $\eta$ and
$\rho$ respectively describe the fractional power and the contribution
of the common environment to the path variance. The exponent is an
independent scalar object, whose existence and required uniform
estimates are proved in this paper. The factors $P_e$ retain this
conditional information. Their product is established by a compatible
construction and successive conditional estimates; it does not require
independence of all the underlying path events.

This description builds on the order-statistics and clustering ideas
of Ford and Koukoulopoulos. Fractional moments themselves already
appear in \cite[Section~1]{KoukLocal}. The additional issue here is to
realize all the constraints simultaneously while preserving the common
environment, the original coordinates and the powers attached to
different owners. These requirements become essential in the
high-dimensional configurations discussed in \cite[Section~2.2]{Kouk}.
The definition of $\Phi_k$ extends the positive-length expression by a
specified common lift of every original coordinate and an infimum over
that lift. The saddle and all its associated data are recomputed at each
lift. An arithmetic stability argument proves that this definition also
covers zero and bounded original lengths.

\subsection{The proof}

We use Koukoulopoulos's arithmetic reduction
\cite[Theorem~1.7]{Kouk} to express the normalized divisor count as a
harmonic sum of volumes of unions of divisor boxes. The main analytic
statement evaluates this sum by $\Phi_k$. Its lower bound is obtained
by realizing the prescribed constraints on the original arithmetic
space, retaining the same accepting density through the successive
probabilistic and arithmetic comparisons. The upper bound uses causal
coverings and local conditional estimates to show that the same losses
are necessary. Finally, lift removal and the multiplication-table
transfer give the two theorems above.

The probability arguments include the scalar persistence theorem,
finite bridges with exact totals, and the simultaneous treatment of
several comparisons. The ordinary order-statistics estimates enter
separately. For the ordered-sum upper estimate, the method of
\cite[Lemma~5.3]{KoukLocal} requires an integer threshold in its band
decomposition. We give the corrected argument in full, starting from
its Lemmas~5.1 and~5.2, and deduce the real-threshold, all-positive-count
case needed here. The ordinary lower estimate is the separate
\cite[Lemma~8.3]{Kouk}.
All proofs needed for the stated results, including the boundary
transfers, are contained in this paper and its technical appendices.
The numbering of results cited from \cite{Kouk} refers throughout to
arXiv:1102.3236v3; for \cite{KoukLocal} it refers to the author version
linked in the bibliography.
